\section{Recovering \texorpdfstring{$f_6$}{f6} from six-qubit AME signatures}
\label{sec:recover-f6}

In this section, following the~\cref{sec: Third Order Orth}'s proof, we have realized a 6-ary signature $f$ and its normalized state is an $\mathrm{AME}(6,2)$ state. Then we can either get $\Holant(\F)$ is $\#\operatorname{P}$-hard, or realize an special 6-ary signature $f_6$ and four Bell signatures $\mathcal{B}$ after a holographic transformation by a real orthogonal matrix $O$, i.e., $\Holant(O\F,f_6,\mathcal{B})\leq_T \Holant(\F)$. With this $f_6$ and the four Bell signatures $\mathcal{B}$ in hand, in the next~\cref{sec: f6 is hard}, we will prove that $\Holant(O\mathcal{F},f_6,\mathcal{B})$ is $\#\operatorname{P}$-hard. 

We retain the standing assumptions that $\mathcal F$ is closed under entrywise conjugation, and is
outside condition \textup{(T)}.

Now we give the definition of the 6-ary signature $f_6$,
\[
 f_6(x)=\mathbf1_{x_1+\cdots+x_6=0}
 (-1)^{x_1+x_2+x_3+q(x)},\qquad
 q(x)=x_1x_2+x_1x_3+x_2x_3+x_1x_4+x_2x_5+x_3x_6,
\]
where Boolean arithmetic is over $\mathbb{F}_2$. Its normalized state is an $\mathrm{AME}(6,2)$ state. Now we have a 6-ary signature $f$ and its normalized state is $\mathrm{AME}(6,2)$, since $\mathrm{AME}(6,2)$ has one local-unitary normal form (see~\cref{sec: AME62 has a common local unitary normal} for the detailed proofs), after relabeling its inputs, we can get that 
\begin{equation}\label{eq:f6-local-unitary-form}
    f=\lambda(U_1\otimes\cdots \otimes U_6)f_6, \quad \lambda\neq 0, \quad U_iU^\dagger=I.
\end{equation}
However, the matrices $U_1,\dots,U_6$ maybe not available in $\Holant(\F)$, so instead of proving that we can directly realize $f_6$ in $\Holant(\F)$,
we will prove that we can convert $U_1,\dots,U_6$ to $V_1,\dots,V_6$ where $[V_i]\in \Gamma$ for $1\leq i\leq 6$ after a orthogonal holographic transformation. Let $h=O^{\otimes6}f=\lambda(V_1\otimes\cdots\otimes V_6)f_6$.
For each pair of variables, we contract $h$ using $=_2$ and extract
the two resulting binary signatures. Let $G$ be the group generated by
the projective classes of these factors, their conjugates, and their
transposes. We then treat the cases $G=\Gamma$ and $G\subseteq K_4$
to remove the remaining local matrices and obtain $f_6$ together
with the four Bell signatures.

The main theorem of the section is the following theorem.


\begin{theorem}[Realizing $f_6$]
\label{thm:f6-equality-realization}
Suppose that $\mathcal F$ contains a nonzero six-ary signature whose
normalized state is $\operatorname{AME}(6,2)$. Then either
$\operatorname{Holant}(\mathcal F)$ is $\#\mathrm P$-hard, or there
is a real orthogonal matrix $O$ such that
\[
 \operatorname{Holant}(O\mathcal F\cup\{f_6\}\cup\mathcal B)
 \leq_T \operatorname{Holant}(\mathcal F).
\]
\end{theorem}






\subsection{Binary contractions and hardness}
\label{subsec:f6-binary-contractions}

In this subsection, we first analyze binary contractions of $f_6$,
then use the local-unitary expression of $f$ to constrain the matrices
$U_i$. By the preceding reductions and \cref{lm: base case 2n=4},
realizing a nonzero four-ary signature outside $\mathcal U^\otimes$
gives hardness. Thus, unless hardness already follows, every nonzero
contraction of $f$ through a realizable binary signature belongs to
$\mathcal U^\otimes$.

Recall that the matrices $I=\left(\begin{matrix}
    1 & 0\\
    0 & 1\\
\end{matrix}\right),X=\left(\begin{matrix}0&1\\1&0\end{matrix}\right)$,
$Z=\left(\begin{matrix}1&0\\0&-1\end{matrix}\right)$, and
$J=\left(\begin{matrix}0&1\\-1&0\end{matrix}\right)$. For a nonzero matrix $A$, $[A]$ denotes its class up to a nonzero scalar. Also recall the following notation of Kelin group and tetrahedral,
\[
 C_{\epsilon_1,\epsilon_2,\epsilon_3}
 =\tfrac12(I+\epsilon_1\mathrm iX+\epsilon_2J+\epsilon_3\mathrm iZ),
 \qquad
 \Gamma=K_4\cup
 \{[C_{\epsilon_1,\epsilon_2,\epsilon_3}]:
          \epsilon_1,\epsilon_2,\epsilon_3\in\{\pm1\}\},
\]
where $K_4=\{[I],[X],[Z],[J]\}$. $K_4$ is the projective Kelin group and $\Gamma$ is the projective tetrahedral group. Both are closed under multiplication, inverse, transpose, and conjugation.
The matrices in $\mathcal B=\{I,X,Z,J\}$ are the four Bell signatures.

Also recall the gadget construction of $\Holant(\F)$, for a 6-ary signature $h$ and a binary matrix $Q$, write $\partial_{ij}^Qh=\sum_{x_i,x_j}Q(x_i,x_j)h(x)$, which is called a contraction of $h$ through $Q$ on variables $x_i,x_j$. In particular,
$\partial_{ij}^Ig$ joins ports $i,j$ directly by $=_2$.
We first determine when such a contraction of $f_6$ splits into two
binary factors.

\begin{lemma}[Binary contractions of $f_6$]
\label{lem:f6-contraction-test}
For every nonzero binary matrix $Q$ and every pair $1\le i<j\le6$,
the signature $\partial_{ij}^Qf_6$ is nonzero. It is a tensor product
of two binary signatures if and only if $[Q]\in\Gamma$.
In that case, both factors are projectively unitary and their
projective classes belong to $\Gamma$.
\end{lemma}
\begin{proof}[Proof of Lemma~\ref{lem:f6-contraction-test}]
Without loss of generality, we consider the first two variables $x_1,x_2$.
First take $i=1,j=2$ and write
$Q=\left(\begin{smallmatrix}a&b\\c&d\end{smallmatrix}\right)$.
By the definition of $f_6$, with row and column indices ordered
as $00,01,10,11$, we have
\[
M_{x_3x_4,x_5x_6}(\partial_{12}^Qf_6)=
\begin{pmatrix}
a-d&-b-c&b-c&a+d\\
c-b&a+d&a-d&b+c\\
-b-c&a-d&-a-d&c-b\\
-a-d&b-c&b+c&a-d
\end{pmatrix}.
\]
This matrix is nonzero when $Q\ne0$: its entries determine
$a-d,a+d,b-c,b+c$ and hence $a,b,c,d$.
The equality $\partial_{12}^Qf_6=A(x_3,x_4)B(x_5,x_6)$ holds
for nonzero binary signatures $A,B$ if and only if this matrix
has rank one. The other two possible factorizations are described
by $M_{x_3x_5,x_4x_6}(\partial_{12}^Qf_6)$ and
$M_{x_3x_6,x_4x_5}(\partial_{12}^Qf_6)$ in the same way.
A nonzero matrix has rank one if and only if all its
$2$-by-$2$ minors vanish. Applying this criterion gives the following
necessary conditions for each matrix:
\[
\begin{array}{c|l}
M_{x_3x_4,x_5x_6}&a^2+d^2=ab-cd=b^2+d^2=ac+bd=ad+bc=c^2-d^2=0,\\
M_{x_3x_5,x_4x_6}&a^2+d^2=ab+cd=b^2-d^2=ac-bd=ad+bc=c^2+d^2=0,\\
M_{x_3x_6,x_4x_5}&a^2-d^2=ab=b^2-c^2=ac=bd=cd=0.
\end{array}
\]
Here all three matrices are evaluated at $\partial_{12}^Qf_6$.
For example, in the displayed matrix the minors on rows $1,4$
and columns $1,4$, and on rows $1,2$ and columns $2,3$, are
$2(a^2+d^2)$ and $-2(ab-cd)$, respectively.
The last system gives either $b=c=0$, $a=\pm d\ne0$, or
$a=d=0$, $b=\pm c\ne0$, hence exactly the classes in $K_4$.
In either of the first two systems, $d=0$ forces $Q=0$.
Scaling to $d=1$, their solutions are respectively
$(a,b,c)=(\epsilon\mathrm i,\eta\mathrm i,-\epsilon\eta)$ and
$(\epsilon\mathrm i,\eta,-\epsilon\eta\mathrm i)$, where
$\epsilon,\eta\in\{\pm1\}$.
Their classes are $[C_{\eta,\epsilon\eta,\epsilon}]$ and
$[C_{-\epsilon\eta,\eta,\epsilon}]$, giving the remaining eight
classes of $\Gamma$.

Substituting these solutions gives the following factorizations,
up to a nonzero scalar. In a row with variable pairs $(x_a,x_b)$
and $(x_c,x_d)$, the last column lists the classes of
$A,B$ in $\partial_{12}^Qf_6=A(x_a,x_b)B(x_c,x_d)$.
\[
\begin{array}{c|c|c}
[Q]&\text{variable pairs}&\text{factor classes}\\ \hline
{[I]}&(x_3,x_6),(x_4,x_5)&[J],\ [X]\\
{[X]}&(x_3,x_6),(x_4,x_5)&[X],\ [Z]\\
{[Z]}&(x_3,x_6),(x_4,x_5)&[I],\ [I]\\
{[J]}&(x_3,x_6),(x_4,x_5)&[Z],\ [J]\\
{[C_{\eta,\epsilon\eta,\epsilon}]}&(x_3,x_4),(x_5,x_6)&
[ZC_{\eta,-\epsilon\eta,\epsilon}],\ [C_{-\eta,-\epsilon\eta,\epsilon}]\\
{[C_{-\epsilon\eta,\eta,\epsilon}]}&(x_3,x_5),(x_4,x_6)&
[ZC_{-\epsilon\eta,\eta,\epsilon}],\ [C_{\epsilon\eta,\eta,\epsilon}]
\end{array}
\]
Every listed class lies in $\Gamma$ and has a unitary representative.
This proves the converse as well. For each solution only one of the
three systems holds, and the two factors for that partition are
unique up to reciprocal scalars. Thus the conclusion holds for
any choice of factors.

\end{proof}

Contracting a binary signature $B$ with inputs $i,j$ of $f$ corresponds
to contracting $U_i^TBU_j$ with the same inputs of $f_6$.
The following theorem gives hardness when this relative matrix lies
outside $\Gamma$. We first use its equality case, $B=I$, to obtain
the conditions needed for a common change of basis.

\begin{theorem}[The relation between $U_i$ and $U_j$]
\label{thm:f6-nongamma-hard}
Let $f$ be the realized signature in \eqref{eq:f6-local-unitary-form}.
Suppose that a nonzero binary signature $B$ is available in $\Holant(\F)$.
If $[U_i^TBU_j]\notin\Gamma$ for some $i<j$, then
$\Holant(\mathcal F)$ is $\#\mathrm P$-hard.
In particular, hardness follows if $[U_i^TU_j]\notin\Gamma$ for some
$i<j$.
\end{theorem}

\begin{proof}[Proof of Theorem~\ref{thm:f6-nongamma-hard}]
We use the standing assumptions that $\mathcal F$ is fixed, finite,
algebraic, conjugate closed, and outside \textup{(T)}.
If $B$ is realized by a gadget, conjugating its vertex labels
realizes $\bar B$. More generally, the same proof applies when
$B,\bar B$ can be jointly adjoined by a polynomial-time Turing
reduction: work over that enlarged conjugate-closed family and
compose the reductions at the end.
Contract $B$ with inputs $i,j$ of $f$. The resulting signature is
\[
 g=\partial_{ij}^Bf
 =\lambda\left(\bigotimes_{k\ne i,j}U_k\right)
       \partial_{ij}^{\,U_i^TBU_j}f_6.
\]
The matrix on the right is $U_i^TBU_j$ because its $(a,b)$ entry is
$\sum_{s,t}U_i(s,a)B(s,t)U_j(t,b)$.
If its class lies outside $\Gamma$, \cref{lem:f6-contraction-test}
shows that $g$ is nonzero and is not a product of two binary signatures;
the remaining local unitaries preserve these properties.
In particular, $g\notin\mathcal U^\otimes$, since every nonzero
four-ary signature in $\mathcal U^\otimes$ is a product of two
binary signatures. By \cref{lm: base case 2n=4}, this gives hardness.
If $B,\bar B$ were adjoined, the enlarged family still contains
$\mathcal F$ and remains outside \textup{(T)}; composing the
reductions gives hardness of $\Holant(\mathcal F)$.
Taking $B=I$ proves the last assertion, since this contraction
simply joins the two inputs by an equality edge.
\end{proof}

Since $=_2$ is available, it remains to consider the case
$[U_i^TU_j]\in\Gamma$ for every $i<j$.


\subsection{A common real orthogonal transformation}
\label{subsec:f6-real-orthogonal}

Unless $\Holant(\F)$ is $\#\operatorname{P}$-hard, \cref{thm:f6-nongamma-hard} gives
$[U_i^TU_j]\in\Gamma$ for every $i<j$. The next lemma gives one real
orthogonal holographic transformation that converts all six local matrices $U_i$ into
$V_i$ which belongs to $\Gamma$ up to a scalar.

\begin{lemma}[A common orthogonal transformation]
\label{lem:f6-common-orthogonal}
Let $U_1,\ldots,U_6$ be unitary matrices such that
$[U_i^TU_j]\in\Gamma$ for every $i<j$.
Then there are a real orthogonal matrix $O$ and unitary matrices
$V_1,\ldots,V_6$ such that $OU_i=V_i$ and $[V_i]\in\Gamma$ for
every $i$. 
\end{lemma}

\begin{proof}[Proof of Lemma~\ref{lem:f6-common-orthogonal}]
Put $A_i=U_1^TU_i$ for $1\le i\le6$. These matrices are unitary,
and $[A_i]\in\Gamma$ for $i>1$. Since
$U_2^TU_3=A_2^TA_1^{-1}A_3$, closure of $\Gamma$ gives
$[A_1^{-1}]=[A_2^{-T}(U_2^TU_3)A_3^{-1}]\in\Gamma$ as well.

Let $R=\overline{U_1}$. Then $U_i=RA_i$ and
$S=R^TR=A_1^{-1}$ is symmetric and unitary.
The only symmetric classes in $\Gamma$ are $[I],[X],[Z]$:
$J$ is antisymmetric, and the off-diagonal entries of
$C_{\epsilon_1,\epsilon_2,\epsilon_3}$ differ by $\epsilon_2$.
Multiplying $R$ by a phase and all $A_i$ by its inverse, we may
assume $S\in\{I,X,Z\}$.

The following choices are unitary, have classes in $\Gamma$, and
satisfy $A_0^TSA_0=I$:
\[
\begin{array}{c|ccc}
S&I&X&Z\\ \hline
A_0&I&e^{\mathrm i\pi/4}C_{+,+,+}^{-1}
       &e^{-\mathrm i\pi/4}C_{+,+,-}^{-1}.
\end{array}
\]
Indeed, $C_{+,+,+}^TC_{+,+,+}=\mathrm iX$ and
$C_{+,+,-}^TC_{+,+,-}=-\mathrm iZ$.
Hence $RA_0$ is both unitary and orthogonal, so it is real.
Set $O=(RA_0)^T$ and $V_i=A_0^{-1}A_i$.
Then $O$ is real orthogonal, $OU_i=V_i$, and each $V_i$ is unitary
with $[V_i]\in\Gamma$.
\end{proof}

We now can realize
\begin{equation}
 h:=O^{\otimes6}f=\lambda(V_1\otimes\cdots\otimes V_6)f_6,
 \label{eq:f6-transformed-signature}
\end{equation}
on the $\Holant(O\mathcal{F})$, and 
$\Holant(\mathcal F)\equiv_T\Holant(\mathcal F)$.

In this basis, $[V_i^TBV_j]\in\Gamma$ if and only if $[B]\in\Gamma$,
since $[V_i],[V_j]\in\Gamma$ and $\Gamma$ is closed under transpose,
multiplication, and inversion. We can therefore state the hardness
condition directly for a binary signature realized over $\mathcal E$.

\begin{lemma}[Non-$\Gamma$ hard]
\label{lem:f6-nongamma-binary}
Under the standing assumptions, let $\mathcal E=O\mathcal F$ and $h$
be as above. If $\mathcal E$ realizes a nonzero binary signature $B$
with $[B]\notin\Gamma$, then $\Holant(\mathcal E)$, and hence
$\Holant(\mathcal F)$, is $\#\mathrm P$-hard.
\end{lemma}

\begin{proof}[Proof of Lemma~\ref{lem:f6-nongamma-binary}]
Since $[V_1],[V_2]\in\Gamma$ and $\Gamma$ is closed under transpose,
multiplication, and inversion, $[B]\notin\Gamma$ implies
$[V_1^TBV_2]\notin\Gamma$.
Apply \cref{thm:f6-nongamma-hard} to $h$ over $\mathcal E$,
this gives that $\#\operatorname{P}$-hardness of $\Holant(\mathcal{E})$.
\end{proof}

It remains to recover $f_6$ from the available signature $h$.


\subsection{Recovering \texorpdfstring{$f_6$}{f6} from binary factors}
\label{subsec:f6-binary-factors}

We define the projective binary group of $\Holant(O\mathcal{F})$, denoted by
\begin{equation}
 G:=\{[B]: B\ne 0 \text{ is a binary signature that can be realized in }\Holant(O\mathcal{F}).\}.
\end{equation}
where the $=_2$ is included, so $[I]\in G$.
By \cref{lem:f6-nongamma-binary}, either
$\Holant(O\mathcal F)$ is $\#\mathrm P$-hard or $G\subseteq\Gamma$.
We henceforth consider the latter case.

Joining binary gadgets gives matrix multiplication, exchanging their
inputs gives transposition, and conjugating their vertex labels gives
entrywise conjugation. Hence $G$ is closed under these operations.
Since $\Gamma$ is finite, inverses are powers, so $G$ is a subgroup
of $\Gamma$. In particular, every binary factor of an equality
contraction of $h$ has its class in $G$.

\begin{lemma}[Two possibilities for the binary group]
\label{lem:f6-group-alternatives}
For the projective binary group $G$ of $\Holant(O\F)$, either
$G=\Gamma$ or $G\subseteq K_4$.
\end{lemma}

\begin{proof}
Suppose $G\nsubseteq K_4$, and choose
$[C_{a,b,c}]\in G\setminus K_4$.
Since $G$ is closed under transpose and multiplication, it contains
the class of
\[
C_{a,b,c}^TC_{a,b,c}
=\begin{cases}
a\mathrm iX,&c=ab,\\
c\mathrm iZ,&c=-ab.
\end{cases}
\]
Thus $G$ contains a nonidentity element of $K_4$.
Conjugation by $[C_{a,b,c}]$ cyclically permutes
$[X],[Z],[J]$, so $K_4\subseteq G$.
The image of $[C_{a,b,c}]$ generates $\Gamma/K_4$, which has
order three. Therefore $G=\Gamma$.
\end{proof}

We handle these two cases separately.
If $G=\Gamma$, every local inverse $[V_i^{-1}]$ is realizable
over $\mathcal E'$, so the local matrices can be removed directly.
If $G\subseteq K_4$, all binary factors of the equality contractions
of $h$ lie in $K_4$. We use this restriction to analyze $h$;
it does not by itself imply that each $[V_i]$ lies in $K_4$.

\begin{lemma}[The tetrahedral case]
\label{lem:f6-tetrahedral-case}
For the projective binary group $G$ of $\Holant(O\F)$,
if $G=\Gamma$, then
\[
\Holant(O\mathcal F\cup\{f_6\}\cup\mathcal B)
\leq_T\Holant(O\mathcal F).
\]
\end{lemma}
\begin{proof}[Proof of Lemma~\ref{lem:f6-tetrahedral-case}]
Recall that
$h=\lambda(V_1\otimes\cdots\otimes V_6)f_6$, where
$\lambda\ne0$ and $[V_i]\in\Gamma$.
Since $G=\Gamma$, each $[V_i^{-1}]$ is available.
Attaching these binary signatures to the corresponding inputs of $h$
gives, up to a nonzero scalar,
$
(V_1^{-1}\otimes\cdots\otimes V_6^{-1})h=\lambda f_6.
$
The four Bell classes also belong to $G$, so all signatures in
$\mathcal B$ are available.
Thus,
$
\Holant(O\mathcal F\cup\{f_6\}\cup\mathcal B)
\leq_T\Holant(O\mathcal F).
$
\end{proof}

\begin{lemma}[The Klein case]
\label{lem:f6-klein-case}
For the projective binary group $G$ of $\Holant(O\F)$,
if $G\subseteq K_4$, then $G=K_4$ and 
\[
\Holant(O\mathcal F\cup\{f_6\}\cup\mathcal B)
\leq_T\Holant(O\mathcal F).
\]
\end{lemma}


\begin{proof}[Proof of Lemma~\ref{lem:f6-klein-case}]
Every binary factor of an equality contraction of $h$ has its
projective class in $G\subseteq K_4$.
We first use this restriction to show that, after relabeling its
inputs, $h$ differs from $f_6$ only by Bell matrices.
We then obtain the Bell signatures and remove these matrices.

For the calculation, use the fixed representative
$g=(I\otimes I\otimes(-J)\otimes X\otimes I\otimes I)f_6$.
The definition of $f_6$ gives
\[
\begin{aligned}
g(x)={}&I(x_1,x_2)I(x_3,x_6)I(x_4,x_5)\\
&+(\mathrm iX)(x_1,x_2)(\mathrm iZ)(x_3,x_6)J(x_4,x_5)\\
&+J(x_1,x_2)(\mathrm iX)(x_3,x_6)(\mathrm iZ)(x_4,x_5)\\
&+(\mathrm iZ)(x_1,x_2)J(x_3,x_6)(\mathrm iX)(x_4,x_5).
\end{aligned}
\]
Thus $h=\nu(B_1\otimes\cdots\otimes B_6)g$, where $\nu\ne0$
and $[B_i]\in\Gamma$.
Put $r=C_{+,+,+}$, and let $c_i\in\{0,1,2\}$ specify the
coset $[B_i]\in[r]^{c_i}K_4$.
All calculations with these indices are modulo $3$.
Multiplication adds coset indices, while transposition negates them:
indeed, $C_{a,b,c}^T=C_{a,-b,c}$, and
$[C_{a,b,c}]\in[r]K_4$ exactly when $abc=1$.
Consequently $B_i^TB_j$ has coset index $c_j-c_i$.

The four-term expansion of $g$ gives the following contractions.
Each row expresses $\partial_{ij}^Qg$, up to a nonzero scalar,
as $A(x_a,x_b)B(x_c,x_d)$.
The last column gives the coset indices of $[A],[B]$.
\[
\begin{array}{cc|c|c|c}
ij&Q&\text{variable pairs}&(A,B)&\text{coset indices}\\ \hline
12&I&(36)(45)&(I,I)&(0,0)\\
12&r&(34)(56)&
\left(
\begin{pmatrix}1&\mathrm i\\-1&\mathrm i\end{pmatrix},
\begin{pmatrix}1&-1\\-\mathrm i&-\mathrm i\end{pmatrix}
\right)&(2,1)\\
12&r^2&(35)(46)&
\left(
\begin{pmatrix}1&\mathrm i\\1&-\mathrm i\end{pmatrix},
\begin{pmatrix}-1&-1\\\mathrm i&-\mathrm i\end{pmatrix}
\right)&(2,1)\\
13&I&(25)(46)&(Z,Z)&(0,0)\\
13&r&(24)(56)&
\left(
\begin{pmatrix}1&\mathrm i\\-1&\mathrm i\end{pmatrix},
\begin{pmatrix}1&\mathrm i\\1&-\mathrm i\end{pmatrix}
\right)&(2,2)\\
13&r^2&(26)(45)&
\left(
\begin{pmatrix}1&-\mathrm i\\1&\mathrm i\end{pmatrix},
\begin{pmatrix}-1&1\\\mathrm i&\mathrm i\end{pmatrix}
\right)&(2,1).
\end{array}
\]
For example, contracting inputs $1,2$ replaces the first binary
matrix in each term of the expansion by its entrywise contraction
with $Q$. Substituting $Q=I,r,r^2$ and factoring gives the first
three rows. Summing over inputs $1,3$ gives the remaining rows.

The variable pairs and coset indices in each row depend only on
the coset of $[Q]$.
To see this, the same expansion shows that
\[
X_1X_4Z_5Z_6,\qquad
X_2X_4Z_3Z_5,\qquad
X_3X_4Z_2Z_6,\qquad
X_4X_5Z_1Z_2
\]
fix $g$, where a subscript indicates the input on which a matrix
acts and all other inputs carry $I$.
Their restrictions to inputs $1,2$ are
$X_1,X_2,Z_2,Z_1Z_2$, and their restrictions to inputs $1,3$ are
$X_1,Z_3,X_3,Z_1$.
In both cases these restrictions generate all two-input Bell
classes. Hence Bell matrices on the contracted inputs can be
transferred to the remaining inputs. They do not change the
variable pairs or the coset indices of the resulting factors.

Now consider an equality contraction of $h$ at inputs $i,j$.
It corresponds to contracting $B_i^TB_j$ into $g$.
A factor $A$ on remaining inputs $a,b$ becomes $B_aAB_b^T$.
If $[A]\in[r]^uK_4$, this factor has coset index $c_a+u-c_b$,
which must be zero because its class belongs to $G\subseteq K_4$.

These constraints depend only on differences of the $c_i$.
Subtracting $c_1$ from all six indices, the rows for $12$ and $13$
give the following possibilities:
\[
\begin{array}{c|ccc}
 &c_3=0&c_3=1&c_3=2\\ \hline
c_2=0&000000&001221&\text{none}\\
c_2=1&010212&\text{none}&012120\\
c_2=2&\text{none}&021102&022011.
\end{array}
\]
Each entry lists $(c_1,\ldots,c_6)$.
For example, $(c_2,c_3)=(0,1)$ gives
$c_3=c_6$, $c_4=c_5$, $c_2+2=c_4$, and $c_5+2=c_6$,
hence $001221$.
The other eight choices are obtained from the corresponding
rows of the contraction table in the same way.
Restoring the common shift gives at most eighteen tuples.

We next remove these coset indices using identities of $g$.
For $(P_\pi g)(x)=g(x_{\pi(1)},\ldots,x_{\pi(6)})$, the following
rows satisfy
$(S_1\otimes\cdots\otimes S_6)g=\mu P_\pi g$
for some $\mu\ne0$:
\[
\begin{array}{c|c}
(S_1,\ldots,S_6)&(\pi(1),\ldots,\pi(6))\\ \hline
(I,I,I,\mathrm iZ,\mathrm iZ,I)&(5,2,6,1,3,4)\\
(r,r,r,r,r,C_{-,+,-})&(6,5,3,1,4,2)\\
(I,I,r,C_{+,-,+},C_{+,-,+},C_{+,-,-})&(6,4,5,3,2,1).
\end{array}
\]
These identities follow by applying each row to the four-term
expansion of $g$: a binary factor $M$ on inputs $a,b$ becomes
$S_aMS_b^T$, and substitution gives the indicated permutation
of the same tensor, up to scalar.

Composition of these identities permutes and adds their coset
tuples. More explicitly, if
$L_Sg=\mu P_\pi g$ and $L_Tg=\eta P_\sigma g$, where
$L_S=\bigotimes_iS_i$ and $L_T=\bigotimes_iT_i$, then
\[
(P_\pi L_TP_\pi^{-1})L_Sg
=\mu\eta P_\pi P_\sigma g.
\]
The third row has tuple $001221$.
Composing with the first row permutes a tuple to
$(c_4,c_2,c_5,c_6,c_1,c_3)$.
After subtracting its first coordinate, this gives the cycle
\[
001221\longmapsto010212\longmapsto022011
\longmapsto021102\longmapsto012120\longmapsto001221.
\]
The identity gives $000000$, and the second row adds a common
shift of one. Thus these identities supply all eighteen tuples.

Choose an identity with $[S_i]K_4=[B_i]K_4$ for every $i$.
Then $D_i=B_iS_i^{-1}$ has class in $K_4$, and
\[
h=\nu\mu(D_1\otimes\cdots\otimes D_6)P_\pi g.
\]
This is a new expression for the same available signature $h$;
it does not require the matrices $S_i$ to be available.
Relabeling the inputs therefore leaves an available signature
equal, up to scalar, to $g$ with Bell matrices on its inputs.

It remains to obtain the Bell signatures.
The expansion of $g$ also gives
\[
g(x)=\mathbf1_{\sum_i x_i=0}(-1)^{q_0(x)+x_6},
\qquad
q_0(x)=x_2x_3+x_3x_6+x_1x_6+x_1x_5+x_2x_5.
\]
A Bell matrix flips an input bit and adds a linear phase,
up to scalar. Thus the relabeled $h$ has the form
$\alpha\mathbf1_{\sum_i x_i=s}
(-1)^{q_0(x)+\sum_i\ell_i x_i}$,
where $\alpha\ne0$ and $s,\ell_i\in\mathbb F_2$.

The following are six binary factors of its equality contractions.
The last two columns identify each factor's class as $[X^uZ^v]$.
\[
\begin{array}{c|c|cc}
\text{contracted inputs}&\text{factor inputs}&u&v\\ \hline
1,2&(x_4,x_5)&s+\ell_1+\ell_2&\ell_4+\ell_5\\
1,3&(x_2,x_5)&\ell_1+\ell_3&
1+\ell_1+\ell_2+\ell_3+\ell_5\\
2,3&(x_1,x_4)&s+1+\ell_2+\ell_3&
1+\ell_1+\ell_2+\ell_3+\ell_4\\ \hline
1,2&(x_3,x_6)&\ell_1+\ell_2&
1+\ell_1+\ell_2+\ell_3+\ell_6\\
1,4&(x_2,x_3)&s+\ell_1+\ell_4&
1+s+\ell_2+\ell_3\\
2,4&(x_1,x_6)&s+\ell_2+\ell_4&
1+s+\ell_1+\ell_6.
\end{array}
\]
To obtain these entries, set $x_i=x_j=t$ and sum over $t$.
For $ij=12,13,23,14,24$, the coefficients of $t$ in the exponent
are respectively
\[
\begin{gathered}
x_3+x_6+\ell_1+\ell_2,\qquad
x_2+x_5+\ell_1+\ell_3,\qquad
1+x_5+x_6+\ell_2+\ell_3,\\
x_5+x_6+\ell_1+\ell_4,\qquad
x_3+x_5+\ell_2+\ell_4.
\end{gathered}
\]
The sum forces this coefficient to vanish.
Together with the parity equation, this gives the two binary
supports; substitution gives their phases.
For example, $ij=12$ gives
$x_6=x_3+\ell_1+\ell_2$ and
$x_5=x_4+s+\ell_1+\ell_2$.
The remaining exponent is, up to a constant,
$(1+\ell_1+\ell_2+\ell_3+\ell_6)x_3
+(\ell_4+\ell_5)x_4$,
which gives both $12$ rows.

All six factor classes belong to $G$.
Projective multiplication of $X^uZ^v$ adds the pairs $(u,v)$
over $\mathbb F_2$.
The first three rows sum to $(1,0)$, and the last three sum to
$(0,1)$. Joining the corresponding factors therefore gives
$[X],[Z]\in G$.
Hence $K_4=\langle[X],[Z]\rangle\subseteq G$, and the assumption
$G\subseteq K_4$ yields $G=K_4$.

All Bell signatures are now available.
We can remove the six Bell matrices from the relabeled $h$ to
recover $g$, and then use the fixed Bell transformation defining
$g$ to recover $f_6$.
Only finitely many fixed constructions are used.
Removing their nonzero scalar factors gives
\[
\Holant(O\mathcal F\cup\{f_6\}\cup\mathcal B)
\leq_T\Holant(O\mathcal F).
\]
\end{proof}


\begin{proof}[Proof of \cref{thm:f6-equality-realization}]
Suppose that $\Holant(\mathcal F)$
is not $\#\mathrm P$-hard.
By the AME normal form, after relabeling inputs, write
$f=\lambda(U_1\otimes\cdots\otimes U_6)f_6$ with unitary $U_i$.
Applying \cref{thm:f6-nongamma-hard} with $B=I$ gives
$[U_i^TU_j]\in\Gamma$ for every $i<j$.
By \cref{lem:f6-common-orthogonal}, there is a real orthogonal
matrix $O$ such that
$h=O^{\otimes6}f=\lambda(V_1\otimes\cdots\otimes V_6)f_6$
with $[V_i]\in\Gamma$.

Let $G$ be the projective binary group of $\Holant(O\mathcal F)$.
By \cref{lem:f6-nongamma-binary}, $G\leq\Gamma$.
Lemma~\ref{lem:f6-group-alternatives} gives either $G=\Gamma$
or $G\subseteq K_4$.
Applying \cref{lem:f6-tetrahedral-case} or
\cref{lem:f6-klein-case}, respectively, and using orthogonal
holographic equivalence, we obtain
\[
\Holant(O\mathcal F\cup\{f_6\}\cup\mathcal B)
\leq_T\Holant(O\mathcal F)
\equiv_T\Holant(\mathcal F).
\]
We are done.
\end{proof}

